% GENERATED FILE. Edit canonical lessons, then run npm run generate:books.
% canonical-source-hash: fnv1a64:68c56a0db7dfb37c
% canonical-lessons: MR-C374-prakriya, MR-C374-tappa, MR-C374-sutra, MR-C374-vyakhya, MR-C374-sastra

\chapter{How Things Work}
\label{ch:mr-how-things-work}
\hlchaptermodality{374}

\begin{chapteropening}
\textbf{By the end of this chapter:} \emph{I can talk about processes, phases, formulas and definitions.}

Complete the last lesson of chapter 374: I can talk about processes, phases, formulas and definitions.
\end{chapteropening}

\section[\texorpdfstring{\mr{प्रक्रिया} (prakriyā)}{प्रक्रिया (prakriyā)}]{\mr{प्रक्रिया} (prakriyā) — a process}

\label{lesson:MR-C374-prakriya}



\begin{quote}
Before the new one: say the Marathi for an analysis, then the Marathi for a method.
\end{quote}

\subsection*{You'll want to know: \mr{प्रक्रिया}}
\textbf{\mr{प्रक्रिया}} — \emph{prakriyā} — \textquotedblleft{}a process\textquotedblright{}.

\begin{quote}
\mr{प्रवेशाची} \mr{प्रक्रिया} \mr{एका} \mr{आठवड्यात} \mr{पूर्ण} \mr{होईल}.
\end{quote}

\begin{quote}
\emph{praveśācī prakriyā ekā āṭhvaḍyāt pūrṇa hoīl.}
\end{quote}

\begin{quote}
The admission process will be complete in a week.
\end{quote}

Feminine.

Feminine: X \mr{ची} \mr{प्रक्रिया}.

\begin{cousinweb}
\textbf{\mr{प्रक्रिया}} (\emph{prakriyā}) is a tatsama, taken from Sanskrit \emph{prakriyā}, \textquotedblleft{}producing; procedure\textquotedblright{} (Monier-Williams).
\end{cousinweb}

\begin{grammarlens}[title={What you've built}]
One more word for giving reasons and opinions.
\end{grammarlens}

\subsection*{Your turn}
\begin{itemize}
\raggedright
  \item \emph{Say it:} \emph{prakriyā}
  \item \emph{Say it:} \emph{prakriyā}, once more
  \item \emph{Say it:} \emph{praveśācī prakriyā ekā āṭhvaḍyāt pūrṇa hoīl.} — The admission process will be complete in a week.
  \item \emph{Recall:} say the Marathi for an analysis, then the Marathi for a method, then say \emph{prakriyā} again
\end{itemize}

\begin{tcolorbox}[breakable,colback=teal!4,colframe=teal!35!black,title={Before you move on}]
What does \mr{प्रक्रिया} mean? (\textquotedblleft{}A process\textquotedblright{}.) Say it once more.
\end{tcolorbox}
\section[\texorpdfstring{\mr{टप्पा} (ṭappā)}{टप्पा (ṭappā)}]{\mr{टप्पा} (ṭappā) — a phase; a stage}

\label{lesson:MR-C374-tappa}



\begin{quote}
Before the new one: say the Marathi for a method, then the Marathi for a process.
\end{quote}

\subsection*{You'll want to know: \mr{टप्पा}}
\textbf{\mr{टप्पा}} — \emph{ṭappā} — \textquotedblleft{}a phase; a stage\textquotedblright{}.

\begin{quote}
\mr{हे} \mr{काम} \mr{तीन} \mr{टप्प्यांत} \mr{होणार} \mr{आहे}.
\end{quote}

\begin{quote}
\emph{he kām tīn ṭappyānt hoṇār āhe.}
\end{quote}

\begin{quote}
This work is going to be done in three phases.
\end{quote}

Masculine. Plural: \mr{टप्पे}, \emph{ṭappe}.

Plural \mr{टप्पे}; oblique \mr{टप्प्या}-.

\begin{cousinweb}
\textbf{\mr{टप्पा}} (\emph{ṭappā}) is a Marathi noun whose origin is not settled.
\end{cousinweb}

\begin{grammarlens}[title={What you've built}]
One more word for giving reasons and opinions.
\end{grammarlens}

\subsection*{Your turn}
\begin{itemize}
\raggedright
  \item \emph{Say it:} \emph{ṭappā}
  \item \emph{Say it:} \emph{ṭappā}, once more
  \item \emph{Say it:} \emph{he kām tīn ṭappyānt hoṇār āhe.} — This work is going to be done in three phases.
  \item \emph{Recall:} say the Marathi for a method, then the Marathi for a process, then say \emph{ṭappā} again
\end{itemize}

\begin{tcolorbox}[breakable,colback=teal!4,colframe=teal!35!black,title={Before you move on}]
What does \mr{टप्पा} mean? (\textquotedblleft{}A phase; a stage\textquotedblright{}.) Say it once more.
\end{tcolorbox}
\section[\texorpdfstring{\mr{सूत्र} (sūtra)}{सूत्र (sūtra)}]{\mr{सूत्र} (sūtra) — a formula}

\label{lesson:MR-C374-sutra}



\begin{quote}
Before the new one: say the Marathi for a process, then the Marathi for a phase.
\end{quote}

\subsection*{You'll want to know: \mr{सूत्र}}
\textbf{\mr{सूत्र}} — \emph{sūtra} — \textquotedblleft{}a formula\textquotedblright{}.

\begin{quote}
\mr{गणिताचं} \mr{हे} \mr{सूत्र} \mr{लक्षात} \mr{ठेवायला} \mr{सोपं} \mr{आहे}.
\end{quote}

\begin{quote}
\emph{gaṇitācaṁ he sūtra lakṣāt ṭhevāylā sopaṁ āhe.}
\end{quote}

\begin{quote}
This maths formula is easy to remember.
\end{quote}

Neuter.

Neuter; also a thread, and a short rule in old texts.

\begin{cousinweb}
\textbf{\mr{सूत्र}} (\emph{sūtra}) is a tatsama, taken from Sanskrit \emph{sūtra}, \textquotedblleft{}thread, cord\textquotedblright{} (CDIAL 13561); a short rule is a thread to hold on to. The same entry lists the inherited Marathi \emph{sūt}, \textquotedblleft{}thread\textquotedblright{}.
\end{cousinweb}

\begin{grammarlens}[title={What you've built}]
One more word for giving reasons and opinions.
\end{grammarlens}

\subsection*{Your turn}
\begin{itemize}
\raggedright
  \item \emph{Say it:} \emph{sūtra}
  \item \emph{Say it:} \emph{sūtra}, once more
  \item \emph{Say it:} \emph{gaṇitācaṁ he sūtra lakṣāt ṭhevāylā sopaṁ āhe.} — This maths formula is easy to remember.
  \item \emph{Recall:} say the Marathi for a process, then the Marathi for a phase, then say \emph{sūtra} again
\end{itemize}

\begin{tcolorbox}[breakable,colback=teal!4,colframe=teal!35!black,title={Before you move on}]
What does \mr{सूत्र} mean? (\textquotedblleft{}A formula\textquotedblright{}.) Say it once more.
\end{tcolorbox}
\section[\texorpdfstring{\mr{व्याख्या} (vyākhyā)}{व्याख्या (vyākhyā)}]{\mr{व्याख्या} (vyākhyā) — a definition}

\label{lesson:MR-C374-vyakhya}



\begin{quote}
Before the new one: say the Marathi for a phase, then the Marathi for a formula.
\end{quote}

\subsection*{You'll want to know: \mr{व्याख्या}}
\textbf{\mr{व्याख्या}} — \emph{vyākhyā} — \textquotedblleft{}a definition\textquotedblright{}.

\begin{quote}
\mr{शब्दकोशात} \mr{या} \mr{शब्दाची} \mr{व्याख्या} \mr{दिली} \mr{आहे}.
\end{quote}

\begin{quote}
\emph{śabdakośāt yā śabdācī vyākhyā dilī āhe.}
\end{quote}

\begin{quote}
The dictionary gives the definition of this word.
\end{quote}

Feminine.

X \mr{ची} \mr{व्याख्या} \mr{करणे}, to define X.

\begin{cousinweb}
\textbf{\mr{व्याख्या}} (\emph{vyākhyā}) is a tatsama, taken from Sanskrit \emph{vyākhyā}, \textquotedblleft{}explanation, exposition\textquotedblright{} (Monier-Williams).
\end{cousinweb}

\begin{grammarlens}[title={What you've built}]
One more word for giving reasons and opinions.
\end{grammarlens}

\subsection*{Your turn}
\begin{itemize}
\raggedright
  \item \emph{Say it:} \emph{vyākhyā}
  \item \emph{Say it:} \emph{vyākhyā}, once more
  \item \emph{Say it:} \emph{śabdakośāt yā śabdācī vyākhyā dilī āhe.} — The dictionary gives the definition of this word.
  \item \emph{Recall:} say the Marathi for a phase, then the Marathi for a formula, then say \emph{vyākhyā} again
\end{itemize}

\begin{tcolorbox}[breakable,colback=teal!4,colframe=teal!35!black,title={Before you move on}]
What does \mr{व्याख्या} mean? (\textquotedblleft{}A definition\textquotedblright{}.) Say it once more.
\end{tcolorbox}
\section[\texorpdfstring{\mr{शास्त्र} (śāstra)}{शास्त्र (śāstra)}]{\mr{शास्त्र} (śāstra) — a science; a discipline}

\label{lesson:MR-C374-sastra}



\begin{quote}
Before the new one: say the Marathi for a formula, then the Marathi for a definition.
\end{quote}

\subsection*{You'll want to know: \mr{शास्त्र}}
\textbf{\mr{शास्त्र}} — \emph{śāstra} — \textquotedblleft{}a science; a discipline\textquotedblright{}.

\begin{quote}
\mr{मला} \mr{लहानपणापासून} \mr{शास्त्र} \mr{आवडतं}.
\end{quote}

\begin{quote}
\emph{malā lahānapṇāpāsūn śāstra āvaḍtaṁ.}
\end{quote}

\begin{quote}
I have liked science since I was a child.
\end{quote}

Neuter.

Also at the end of names of fields: \mr{भाषाशास्त्र}, linguistics.

\begin{cousinweb}
\textbf{\mr{शास्त्र}} (\emph{śāstra}) is a tatsama, taken from Sanskrit \emph{śāstra}, \textquotedblleft{}command, advice, treatise\textquotedblright{} (CDIAL 12422).
\end{cousinweb}

\begin{grammarlens}[title={What you've built}]
One more word for giving reasons and opinions.
\end{grammarlens}

\subsection*{Your turn}
\begin{itemize}
\raggedright
  \item \emph{Say it:} \emph{śāstra}
  \item \emph{Say it:} \emph{śāstra}, once more
  \item \emph{Say it:} \emph{malā lahānapṇāpāsūn śāstra āvaḍtaṁ.} — I have liked science since I was a child.
  \item \emph{Recall:} say the Marathi for a formula, then the Marathi for a definition, then say \emph{śāstra} again
\end{itemize}

\begin{tcolorbox}[breakable,colback=teal!4,colframe=teal!35!black,title={Before you move on}]
What does \mr{शास्त्र} mean? (\textquotedblleft{}A science; a discipline\textquotedblright{}.) Say it once more.
\end{tcolorbox}
